\section{Canonical affine sections of a causal array}
\label{sec:sections}
We use finite alphabets and a finite horizon. At a controller stage an
action is chosen before its next report is received. A terminal answer
is an action with a dummy following report. A complete normalized causal
array $C(\mathbf a\mid\mathbf y)$ prescribes the conditional law of all
actions for every externally supplied report word. Nonanticipation means
that the marginal of an action prefix depends only on reports preceding
that prefix. Intermediate buffers that must survive an update are included
as charged cuts. The rows may depend on the declared time, but not on an
unrecorded history, schedule variable or persistent random seed.

At a cut $t$, let $\cP_t$ be the polytope of normalized causal arrays on
the remaining interface. For a controller prefix $h$ whose probability,
with past reports fixed, is positive, let $R_h\in\cP_t$ be its normalized
future array. Remove duplicate rows to form the finite residual set
$\cR_t$. This is a property of $C$, not of a chosen realization. We
require the full array, including unsupported environmental inputs; an
almost-sure partial specification is a different completion problem.
Set
\begin{equation}\label{eq:slice}
 A_t=\aff\cR_t,\qquad L_t=A_t\cap\cP_t,\qquad
 r_t=\dim A_t+1.
\end{equation}
The slice $L_t$ is a compact polytope of dimension $r_t-1$, since it
contains $\cR_t$. Normalization places $A_t$ in a hyperplane not containing
zero, so $r_t$ is also the ordinary row rank of the normalized residuals.

If $R$ is a future array, write $p_R(a)$ for its next-action probability.
It is independent of the next report $y$. For $p_R(a)>0$, denote by
$R^{a,y}$ the normalized continuation after that action and report. For
an actual residual this is another actual residual at the child cut.

\begin{proposition}[Affine slices are closed under positive shifts]
\label{prop:slice-shift}
If $G\in L_t$ and $p_G(a)>0$, then $G^{a,y}\in L_{t+1}$ for every
next report $y$.
\end{proposition}
\begin{proof}
Write $G=\sum_h c_hR_h$ with real coefficients satisfying $\sum_hc_h=1$.
Taking an action marginal gives $p_G(a)=\sum_h c_hp_{R_h}(a)$. Restriction
to $(a,y)$ is linear in the unnormalized arrays, hence
\[
 G^{a,y}=\sum_{h:p_{R_h}(a)>0}
       \frac{c_hp_{R_h}(a)}{p_G(a)}R_h^{a,y}.
\]
The coefficients sum to one. They need not be positive: the display
asserts membership in the affine span at the child, not a convex
realization. Since $G\in\cP_t$ and $p_G(a)>0$, its normalized shift
is a valid nonnegative causal array. Thus it lies in both factors
of \eqref{eq:slice} at the child cut.
\end{proof}

\begin{theorem}[A canonical sufficient criterion and the rank-two case]
\label{thm:slice}
If every $L_t$ is a simplex, a single stochastic clocked machine realizes
$C$ with exactly $r_t$ labels at cut $t$, simultaneously. No realization
can use fewer than $r_t$ at that cut. In particular, every complete
finite causal array with $r_t\le2$ at every cut has a simultaneously
rank-minimal realization. This includes full-support arrays with
arbitrary finite action and report alphabets.
\end{theorem}
\begin{proof}
In any realization, conditioning on a prefix expresses $R_h$ as a convex
combination of the state continuations at cut $t$. Its affine dimension
therefore forces at least $r_t$ labels.

Choose the vertices of $L_t$ as the states at cut $t$. Every actual
residual is their convex combination. From vertex $G$, choose the next
action with probability $p_G(a)$. After the next report, express
$G^{a,y}\in L_{t+1}$ as a convex combination of the child vertices and
use those weights as the update row. An action of probability zero may
have an arbitrary unused update. Proposition~\ref{prop:slice-shift}
ensures that this construction is possible, and the action probability
is independent of the unread report. Initialization uses the
barycentric representation of $C$. Backward induction verifies the entire
specified future law from every vertex state. A simplex of dimension
$r_t-1$ has exactly $r_t$ vertices, proving the claim. A zero-dimensional
slice is a point and a compact one-dimensional slice is a segment, which
proves the last assertion.
\end{proof}

This criterion differs from requiring the residual hull itself to be a
simplex: internal vertices can lie outside that hull. It also differs
from searching for unknown enclosing simplices and unknown shifts
simultaneously. The polytope $L_t$ is fixed by linear equations and
inequalities from the input array; its normalized-shift closure is
proved, not an additional feasibility hypothesis. For fully listed
rational arrays, vertices and barycentric rows can be found by rational
linear algebra and finite polyhedral enumeration. The explicit array
may already be exponentially large in the horizon. No succinct-input
polynomial bound, autonomous row-reuse theorem, or universal equality
between rank and state complexity follows.

The next section shows that the rank-two conclusion is sharp as a
uniform statement about the maximum cut rank. At rank three, separate
rank-optimal enclosures need not be causally compatible even for a
full-support family.
